\scanpage{523}
\mspage{434}
% B01
le plus économique est
% B02
\[ \mathrm{Ob}\,E=\{0\},\quad \mathrm{Fl}(E)=\{\mathrm{id}_0,p\}\quad \text{avec }p^2=p\ne\mathrm{id}_0, \]
% B03
l'enveloppe de Karoubi de $E$ est donnée par
% B04
\diagram{\begin{tikzcd}[ampersand replacement=\&,column sep=small]
 P={}\Bigg\{ \& 1\arrow[r,<->,shift right=2,"\sigma"'] \& 0\arrow[l,shift right=2,"\rho"']\arrow[loop right,"p"] \& \begin{gathered}\rho\sigma=\mathrm{id}_1\\ \sigma\rho=p\ne\mathrm{id}_0\end{gathered}\Big|\Longrightarrow p^2=p\ \Bigg\}.
\end{tikzcd}}

% B05
Cela soulève la question de trouver un énoncé similaire à 4.15.1, donnant des conditions \unsure{servibles|serviables|sensibles} simples, plus fortes que $H'=H(1+2+3)$, \deleted{pour} sur $f:E\longrightarrow C$, pour que $E$ soit accessible et $f$ accessible, en supposant $C$ \uline{karoubienne}. Ou \unsure{encore|mieux}, des conditions néc. et suff. \unsure{servibles|serviables|sensibles} \inserted{(\deleted{\unsure{plus fortes que}}\inserted{plus fortes que} 4.15.4)}, \deleted{\unsure{pour que pour}} sur $f:E\unsure{\relbar|\to} C$, pour que pour tte \uline{petite} catégorie \uline{karoubienne} $C'$, et tout foncteur $C'\relbar C$, la catégorie $E'=E\times_C C'$ soit accessible, et le foncteur $E'\relbar C'$ accessible.

% B06
Une idée naturelle, c'est qu'il suffirait que cette conclusion soit vraie quand $C'$ est \deleted{soit} une des deux catégories $\unsure{\mathbb{\Delta}|\Delta}^1,P$ (où $P$ est la catégorie à \inserted{trois flèches non identiques} décrite plus haut). Mais je doute qu'il en soit ainsi.

% B07
\margin{\unsure{Ça|Ce} devrait marcher par contre si on prend les catégories}
\diagram{\(\begin{tikzcd}[ampersand replacement=\&,column sep=small]2\arrow[r]\&1\arrow[r,shift left=1]\arrow[r,phantom,"\underbrace{\hphantom{1\longrightarrow0\mkern12mu}}_{P}"'{yshift=-1.4ex}]\&0\arrow[l,shift left=1]\arrow[loop right,"p"]\end{tikzcd}\quad\text{et}\quad\begin{tikzcd}[ampersand replacement=\&,column sep=small,row sep=0pt]\&0\arrow[r]\&\unsure{1|{}},\\2\&1\arrow[l]\arrow[r,shift left=1]\arrow[r,phantom,"\underbrace{\hphantom{1\longrightarrow0\mkern12mu}}_{P}"'{yshift=-1.4ex}]\&0\arrow[l,shift left=1]\arrow[loop right,"p"]\end{tikzcd}\ {\unsure{(*)|(+)} }. \)}
\margin{$\nwarrow$ 3 objets et 5 flèches non identiques}

% B08
% slanted
\margin{(\unsure{se précise|équiv.|ce qui donne} : 4.15.\unsure{7|3} \unsure{p.} \deleted{\unsure{0}} \unsure{443|445|435})}

% B09
% slanted
\margin{Ces catégories \deleted{\unsure{représentent}} représentent respectivement les \unsure{deux|{}} foncteurs $\mathrm{Cat}\to\mathrm{Ens}$ qui associent à une catégorie $C$ les diagrammes dans $C$ \deleted{\illegible{2}} formés des \unsure{$y,x$}, \deleted{\unsure{$(y,p,p,u)$}} avec $p$ projecteur, avec \deleted{\unsure{$p,p$}} $p$ \unsure{sur $y$|sur $x$} et $u:y\to x$ (resp. $u:x\to y$), \unsure{(i.e.|(i.e} $u:\deleted{\unsure{y}}\to x$ avec $u=pu$) \unsure{(ou|(resp.} $u:x\to y$ avec \deleted{\unsure{$yp$|$up$}} $u\unsure{=|\cdot}$}

% B10
On va quand même essayer un énoncé un peu très particulier dans ce sens :

% archive sheet number: 522
